% paper/section_2_background.tex
% Full draft — was a stub. Now complete with current literature grounding.

\section{Background}
\label{sec:background}

\subsection{The agent infrastructure attack surface}

Modern LLM agents are composite systems: a language model provides the
reasoning surface, while the harness provides tool dispatch, context
management, multi-turn state, and persistence. In production deployments
the model is typically an external API call, while tool dispatch, context
management, multi-turn state, and persistence are all local harness
code --- a structural asymmetry with a direct security consequence: the
harness, not the model itself, is where most exploitable surface area
lives.

Agent harnesses expose several classes of latent vulnerabilities.
\emph{Configuration-as-code} attacks exploit declarative configuration the
harness executes unsanitized (CVE-2026-21852 on Claude Code's
\texttt{.claude/settings.json}; CVE-2026-21858 on n8n).
\emph{Serialization-injection} attacks exploit tool-response/state
deserialization, enabling arbitrary class instantiation in trusted
namespaces (CVE-2025-68664 on LangChain Core, CVSS 9.3).
\emph{Hook-bypass} attacks exploit incorrectly wired hook points (OpenClaw
issues \#5513, \#5943: \texttt{before\_\allowbreak tool\_\allowbreak call}
documented but never invoked), and \emph{CSRF/session} attacks exploit the
harness's web
surface (CVE-2025-34291 on LangFlow, actively exploited in the wild as of
January 2026~\cite{crowdsec2026langflow}).

These vulnerability classes correspond directly to Kavach's minister
structure: EXECUTOR detects code-execution-class attacks (CVE-2025-59536,
CVE-2025-34291, CVE-2025-68664), VAULT credential-theft-class attacks
(CVE-2026-21852, CVE-2025-68664), CHANNEL exfiltration-class attacks
(CVE-2026-21852 via URL redirection, CVE-2026-21858), and NAVIGATOR
trajectory-manipulation attacks spanning multiple tool calls.
Section~\ref{sec:architecture} gives the MITRE ATT\&CK and OWASP Agentic
2026 taxonomic grounding for each minister.

\subsection{Existing guardrails and their enforcement points}

The guardrail literature can be organized along two axes: (1) \emph{where} in
the agent lifecycle the guardrail enforces, and (2) \emph{what mechanism}
it uses to detect violations.

\paragraph{Input-time guardrails} (PromptGuard 2
\cite{saxe2025llamafirewall}, Azure Prompt Shields) operate on the user's
message before the agent processes it, missing attacks arriving via tool
outputs, the dominant vector in prompt-injection benchmarks
\cite{zhan2024injecagent}.

\paragraph{Plan-time / chain-of-thought guardrails} (AlignmentCheck
\cite{saxe2025llamafirewall}, DRIFT \cite{drift2026}) audit the reasoning
trace before an action commits --- powerful but expensive (full LLM
inference), and blind to \emph{resolved arguments}, seeing only
declarative intent.

\paragraph{Tool-call-time guardrails} operate at the tool-call boundary ---
Kavach's enforcement point. CodeShield \cite{saxe2025llamafirewall} applies
Semgrep to code-returning outputs only; ShieldAgent
\cite{chen2025shieldagent} retrieves rule circuits per trajectory; Task
Shield \cite{jia2025taskshield} checks task alignment. None combines
deterministic parsing of an action's resolved arguments with a semantic
corpus as a secondary triage layer at the tool-call level --- the
combination \S\ref{sec:design} argues for, after finding the semantic
corpus alone was the wrong primary mechanism for three of four detectors.

\paragraph{Information-flow guardrails} (CaMeL \cite{debenedetti2025camel},
FIDES \cite{costa2025fides}, RTBAS \cite{zhong2025rtbas}) tag data and
enforce statically --- provable, but unable to catch attacks whose
provenance is correctly labeled while the argument content is malicious.

\paragraph{Output-time guardrails} (NeMo Guardrails, Conseca
\cite{tsai2025conseca}, and the post-hoc monitor Kavach replaces) act
after execution --- useful for forensics, too late to prevent it.

\paragraph{Drift-detection guardrails} (TaskTracker
\cite{abdelnabi2024tasktracker}, DeepContext
\cite{albrethsen2026deepcontext}, Trajectory Guard
\cite{advani2026trajectoryguard}) measure trajectory divergence from
stated intent across turns. NAVIGATOR sits here at the per-call level,
comparing a proposed consequential action against the session's pinned
intent; we report both a completed taxonomy layer and an incomplete
authorization mechanism in \S\ref{sec:limitations}.

Across these categories, the tool-call boundary is the only enforcement point
where an action's \emph{resolved runtime arguments} are available before it
executes: input-time and plan-time guards see only declared intent, not
concrete arguments; output-time guards observe the action only after
execution; information-flow guards enforce provenance, not argument
semantics. This is Kavach's enforcement point, structurally unreachable by
guards in the other categories. \S\ref{par:external-baselines} reports a
narrower coverage gap for the two closest tool-call-content tools
specifically, not a test of this general claim.

Table~\ref{tab:guardrail-comparison} (\S\ref{sec:related}) places Kavach in
context on five axes: enforcement boundary, mechanism, corpus extensibility,
verdict combination, and reported AgentDojo / InjecAgent performance.

\subsection{Benchmarks and evaluation methodology}

The two agent-security benchmarks we evaluate against are:

\textbf{InjecAgent} \cite{zhan2024injecagent} --- 1,054 prompt-injection
cases (17 user tools $\times$ 62 attacker tools) split into direct-harm (30
tools) and data-stealing (32 tools) settings, each pairing a benign user
instruction with a malicious tool-output payload and an expected attacker
action. Kavach runs the 544-case data-stealing (DS) subset cold, and the
direct-harm subset separately (\S\ref{sec:injecagent}).

\textbf{AgentDojo} \cite{debenedetti2024agentdojo} --- 97 user tasks and
629 security test cases across four environments, scored by ASR and
utility. We evaluate against its structurally applicable ground-truth
attack cases (\S\ref{sec:channel-eval}) and, at scale, against all 629 real
recorded GPT-4o attack trajectories it publishes (\S\ref{sec:gpt4o}).

A methodological note: reporting attack recall without the accompanying
benign false-positive rate conflates detector quality with aggressiveness.
We therefore report both jointly for every detection number in
\S\ref{sec:eval}, following Kapoor et al.~\cite{kapoor2024aiagents}, rather
than gating recall behind a fixed false-positive ceiling.

\subsection{OWASP and MITRE taxonomic grounding}

Kavach's attack-pattern corpus is authored against two widely cited
AI-agent threat taxonomies, not attacks the team happened to see: MITRE
ATT\&CK/ATLAS (EXECUTOR: T1059, T1546; VAULT: T1552, T1539; CHANNEL:
T1041, T1567; NAVIGATOR: T1083 and others) and OWASP Top 10 for Agentic
Applications 2026 (EXECUTOR: ASI04/05; VAULT: ASI03/09; CHANNEL: ASI07/08;
NAVIGATOR: ASI01/10).
